\begin{abstract}
Groundbreaking studies, including BitNet~\cite{bitnet}, are ushering in a transformative era centered on 1-bit Large Language Models (LLMs). In this paper, we present a ternary LLM variant named \textbf{\text{BitNet b1.58}}, wherein every parameter (or weight) adopts a value from the set \{-1, 0, 1\}. This approach delivers performance on par with standard full-precision (such as FP16 or BF16) Transformer-based LLMs with identical architectural size and training data volume, achieving comparable results in perplexity and downstream tasks. Additionally, BitNet b1.58 offers pronounced benefits regarding memory efficiency, reduced latency, improved throughput, and lower energy usage. Of particular significance, this 1.58-bit LLM establishes an innovative scaling law and formulates training methodologies for future LLMs that combine robust performance with enhanced efficiency. The resulting architecture also suggests a novel computational paradigm, paving the way for hardware tailored specifically for 1-bit LLMs.
\end{abstract}

\section{The Era of 1-bit LLMs}

Artificial intelligence has recently witnessed unprecedented advancements through the rapid scaling and capability improvements of Large Language Models (LLMs). Although these systems achieve state-of-the-art outcomes across a variety of natural language processing benchmarks, their escalating parameter counts raise serious deployment and sustainability issues, primarily because of the corresponding spike in energy requirements and operational costs.

A promising strategy to mitigate these drawbacks is post-training quantization, which yields low-bit models optimized for inference~\cite{smoothquant, gptq, quip, quip_sharp}. By lowering the precision of weights and activations, this method enables substantial reductions in computational overhead and memory footprint for LLMs. The progression has advanced from 16-bit models toward even lower-bit implementations, including 4-bit configurations~\cite{gptq, awq}. Nevertheless, despite its widespread adoption in production-level LLMs, post-training quantization remains a fundamentally sub-optimal solution.

Recent advances in 1-bit model structures, exemplified by BitNet~\cite{bitnet}, provide a compelling path toward lowering the operational expenses of LLMs while retaining high accuracy. Standard LLMs operate using 16-bit floating-point representations (specifically FP16 or BF16), with matrix multiplications dominating both memory and computation. As such, the primary computational cost stems from floating-point multiplications and additions. BitNet, in contrast, replaces these costly floating-point operations with integer additions during matrix multiplications, thereby reducing energy requirements for LLMs by several magnitudes. Since many hardware platforms are fundamentally constrained by power budgets, the substantial decrease in energy consumption also leads to potential improvements in computational speed.

Beyond computational aspects, moving model parameters from DRAM to on-chip accelerators' memory (such as SRAM) incurs significant overhead during inference. Although scaling up on-chip SRAM could help throughput, it does so at a considerable financial premium over DRAM. Employing 1-bit LLMs notably shrinks the overall memory footprint—both in terms of bandwidth and total capacity—compared to their full-precision counterparts. This results in more rapid and less costly weight loading from DRAM, translating to greater inference speed and efficiency.

In this study, we propose a novel 1-bit LLM derivative known as \textbf{\text{BitNet b1.58}}, in which each parameter adopts a ternary value from the set \{-1, 0, 1\}. Introducing the zero as an allowable parameter value augments the original BitNet's 1-bit scheme to 1.58 bits under binary encoding. \text{BitNet b1.58} preserves the original model's primary advantages, such as its shift in computational paradigm that nearly eliminates multiplications during matrix operations and is well-suited to hardware optimization. It maintains equivalent energy usage to the foundational 1-bit BitNet, while also providing significant gains over FP16 LLMs in memory usage, data throughput, and latency. Additionally, \text{BitNet b1.58} brings two further benefits: First, by explicitly incorporating zero-valued parameters, the model introduces feature filtering capabilities, thereby enhancing the representational capacity of 1-bit LLMs and boosting their performance. Second, our empirical analysis demonstrates that, beginning at a 3B parameter scale, \text{BitNet b1.58} attains performance on par with FP16 baselines in both perplexity and downstream tasks, given equivalent training and architectural configurations.

\section{BitNet b1.58}

\text{BitNet b1.58} is constructed upon the BitNet architecture, which implements a Transformer model where the conventional \emph{nn.Linear} layers are replaced with \emph{BitLinear} layers. The model is initialized and trained entirely from scratch, employing weights quantized to 1.58 bits and activations quantized to 8 bits. Relative to the initial BitNet formulation, this version incorporates several modifications, which are outlined as follows.

\paragraph{Quantization Function.} We constrain the weights to the set $\{-1, 0, +1\}$ through the use of an \emph{absmean} quantization function. This approach entails normalizing the weight matrix by its mean absolute value and then rounding each resulting element to the nearest integer within $\{-1, 0, +1\}$:

\begin{equation}
    \widetilde{W} = \mathrm{RoundClip}(\frac{W}{\gamma + \epsilon}, -1, 1),
\end{equation}

\begin{equation}
    \mathrm{RoundClip}(x, a, b) = \max(a, \min(b, \mathrm{round}(x))),
\end{equation}

\begin{equation}
    \gamma = \frac{1}{nm}\sum_{ij} |W_{ij}|.
\end{equation}

The procedure for quantizing activations remains identical to that of BitNet, except we omit rescaling activations to the $[0, Q_b]$ interval prior to the application of non-linearities. Instead, activations are rescaled to the interval $[-Q_b, Q_b]$ for each token, thereby eliminating the need for zero-point quantization. This simplification improves implementation and system-level optimization with negligible impact on performance, as observed in our experimental results.

\paragraph{LLaMA-alike Components.}

LLaMA~\cite{llama, llama2} has become the standard architecture for open-source large language models. In alignment with the open-source ecosystem, \text{BitNet b1.58} is designed with LLaMA-inspired components, specifically leveraging RMSNorm~\cite{rmsnorm}, SwiGLU~\cite{swiglu}, and rotary embeddings~\cite{rope}, and omitting all bias terms. These design choices ensure that \text{BitNet b1.58} can be readily incorporated into mainstream open-source frameworks (such as Huggingface, vLLM~\cite{vllm}, and llama.cpp\footnote{\url{https://github.com/ggerganov/llama.cpp}}) with minimal integration overhead.

\section{Results}

We conducted a comparison between \text{BitNet b1.58} and our re-implemented FP16 \text{LLaMA LLM} across multiple model scales. For a consistent evaluation, both model families were pre-trained for 100 billion tokens using the RedPajama dataset~\cite{redpajama}. Zero-shot performance was assessed across a variety of language understanding benchmarks, such as ARC-Easy~\cite{arc}, ARC-Challenge~\cite{arc}, Hellaswag~\cite{hellaswag}, Winogrande~\cite{winoGrande}, PIQA~\cite{piqa}, OpenbookQA~\cite{openbookqa}, and BoolQ~\cite{boolq}. In addition, we present validation perplexity metrics on WikiText2~\cite{wikitext2} and C4~\cite{c4}.

Our analysis also included a detailed assessment of GPU runtime memory usage and inference latency for both \text{LLaMA LLM} and \text{BitNet b1.58}. Measurements were obtained using the FasterTransformer\footnote{\url{https://github.com/NVIDIA/FasterTransformer}} library, which provides highly optimized routines for LLM inference on GPUs. For \text{BitNet b1.58}, the 2-bit kernel provided by Ladder~\cite{ladder} was incorporated. The primary metric reported was time per generated token, which constitutes the main computational expense during inference.

Table~\ref{tab:ppl} presents the perplexity figures along with the associated computational costs for \text{BitNet b1.58} and \text{LLaMA LLM}. The data indicate that, at a scale of 3B parameters, \text{BitNet b1.58} achieves perplexity values that are comparable to the full-precision \text{LLaMA LLM}, while delivering a 2.71$\times$ speedup and reducing GPU memory usage by a factor of 3.55. Notably, \text{BitNet b1.58} with 3.9B parameters is 2.4$\times$ faster, requires 3.32$\times$ less memory, and yields substantially superior results compared to \text{LLaMA LLM} 3B.

A comprehensive breakdown of zero-shot accuracy across downstream tasks is presented in Table~\ref{tab:zero-shot}. For this evaluation, we adopted the procedure from the \emph{lm-evaluation-harness}\footnote{\url{https://github.com/EleutherAI/lm-evaluation-harness}} toolkit. Results highlight that as model capacity grows, the performance disparity between \text{BitNet b1.58} and \text{LLaMA LLM} decreases. More significantly, from the 3B model size onward, \text{BitNet b1.58} achieves parity with the full precision baseline. Consistent with perplexity trends, the end-task evaluations also reveal that \text{BitNet b1.58} 3.9B surpasses \text{LLaMA LLM} 3B, while simultaneously delivering lower memory consumption and reduced latency. These outcomes establish \text{BitNet b1.58} as a Pareto-optimal advance relative to leading LLM architectures.

\paragraph{Memory and Latency}

We extended our analysis by increasing the model sizes to 7B, 13B, and 70B, and systematically assessed the associated computational cost. Figure~\ref{fig:latency-memory-energy} depicts how both latency and memory usage evolve with scaling, indicating that performance gains accelerate as the parameter count rises. Notably, the \text{BitNet b1.58} 70B configuration achieves a speed that is 4.1-fold greater than that of the \text{LLaMA LLM} reference model. This improvement is primarily attributable to the escalating computational demands of the \emph{nn.Linear} component with larger models. Memory usage demonstrates analogous behavior, as the size of the full-precision embedding—remaining constant—becomes a diminishing fraction of total memory as models get larger. Both latency and memory measurements utilized a 2-bit computational kernel, suggesting that there remain opportunities for additional reductions in resource consumption.

\paragraph{Energy}

In addition, we quantified the energy expenditure due to arithmetic operations for both \text{BitNet b1.58} and \text{LLaMA LLM}, with particular attention given to matrix multiplication, which dominates LLM computation costs. Figure~\ref{fig:energy} presents the breakdown of energy sources, revealing that most of \text{BitNet b1.58}'s operations are INT8 additions, whereas \text{LLaMA LLM} relies on both FP16 addition and multiplication. Drawing on the methodology from~\cite{energycost, pokebnn}, our analysis indicates that \text{BitNet b1.58} reduces matrix multiplication energy demands by a factor of 71.4 compared to FP16 operations when implemented on 7nm technology. We further evaluated the total energy usage for processing sequences of 512 tokens. These results demonstrate that energy efficiency gains for \text{BitNet b1.58} become even more pronounced as model scale increases, especially relative to the FP16 \text{LLaMA LLM} baseline. This trend arises because the share of computation in \emph{nn.Linear} operations expands with larger models, while the contribution from other parts of the architecture diminishes.

\paragraph{Throughput}

We assess the throughput capabilities of \text{BitNet b1.58} compared to the \text{LLaMA LLM} with 70B parameters, utilizing two 80GB A100 GPUs and employing pipeline parallelism~\cite{gpipe} to make \text{LLaMA LLM} 70B feasible on this hardware. The batch size was progressively increased up to the maximum that GPU memory could accommodate, while keeping the sequence length fixed at 512. As indicated in Table~\ref{tab:throughput}, the \text{BitNet b1.58} 70B model can accommodate batch sizes as large as 11 times those possible with \text{LLaMA LLM}, thereby achieving a throughput boost of 8.9-fold.

\textbf{\text{BitNet b1.58} introduces a novel scaling dynamic in terms of model efficiency versus inference expense}. Based on the data presented in Figures~\ref{fig:latency-memory-energy} and~\ref{fig:energy}, there is an empirical equivalency between different model scales at 1.58-bit and those at 16-bit precision:

\begin{itemize}
    \item 13B BitNet b1.58 exhibits greater efficiency—with respect to latency, memory consumption, and energy use—than the 3B FP16 LLM.
    \item 30B BitNet b1.58 outperforms the 7B FP16 LLM in terms of latency, memory footprint, and energy requirements.
    \item 70B BitNet b1.58 surpasses the 13B FP16 LLM with superior latency, reduced memory use, and lower energy consumption.
\end{itemize}

\paragraph{Training with 2T Tokens}

The total number of training tokens remains a key variable in large language model development. To examine the token scalability of \text{BitNet b1.58}, we trained a \text{BitNet b1.58} model with 2T tokens, adhering to the dataset composition strategy used by StableLM-3B~\cite{StableLM-3B-4E1T}, a leading 3B-parameter open-source model. Evaluation took place on a benchmark suite composed of Winogrande~\cite{winoGrande}, PIQA~\cite{piqa}, SciQ~\cite{sciq}, LAMBADA~\cite{lambada}, and ARC-easy~\cite{arc}. The zero-shot accuracy is presented in Table~\ref{tab:2t}, where we report the average of the accuracies when both accuracy and normalized accuracy are relevant. Results for StableLM 3b with 2T tokens are cited from its original report. Our experimental evidence demonstrates that \text{BitNet b1.58} consistently outperforms across all assessed tasks, suggesting that 1.58-bit LLMs can achieve robust generalization.

\section{Discussion and Future Work}

\textbf{1-bit Mixture-of-Experts (MoE) LLMs}

Mixture-of-Experts (MoE) architectures have established themselves as an efficient strategy for reducing computational costs in LLMs. However, their widespread adoption has been hampered by substantial memory requirements and heavy inter-device communication. Deploying 1.58-bit LLMs addresses these bottlenecks. The significant drop in memory usage means fewer devices are necessary to run MoE models, while network communication overhead, particularly for activations, is drastically cut down. Ideally, if the memory footprint shrinks sufficiently, these models could potentially be hosted entirely on a single chip, effectively eliminating communication overhead.

\textbf{Native Support for Long Sequences in LLMs}

Managing long sequences is now a foundational requirement for modern LLM applications. One of the primary obstacles in this context is the memory burden caused by key-value (KV) caching. \text{BitNet b1.58} constitutes meaningful progress toward supporting longer context windows natively, as it compresses activations from 16 bits down to 8 bits. This compression permits, for a fixed memory budget, roughly double the sequence length. The potential exists to losslessly compress these representations even further to 4 bits or less with 1.58-bit LLMs; exploring this compression remains an open direction for future research.

\textbf{LLMs for Edge and Mobile Deployment}

Deploying LLMs on edge and mobile platforms remains challenging due to tight constraints on memory and computation. The introduction of 1.58-bit LLMs markedly lowers these requirements, thus enabling their operation on resource-constrained devices where larger, more power-hungry models were impractical. This advances the scope and variety of on-device LLM applications, augmenting device capabilities in novel ways. Furthermore, since most edge and mobile systems rely on CPUs, the efficient execution of 1.58-bit models such as \text{BitNet b1.58} on CPU hardware further bolsters their viability in these scenarios.

\textbf{Innovative Hardware for 1-bit LLMs}

Innovations such as those from Groq\footnote{\url{https://groq.com/}} showcase the promising future of custom hardware, including Linear Processing Units (LPUs), for accelerating LLM workloads. Building on these developments, there is a need and opportunity to design specialized hardware and system architectures targeted specifically at 1-bit LLMs, harnessing the advantages offered by the computation paradigm of BitNet~\cite{bitnet}.